\documentclass{article}
\usepackage[backend=biber,natbib=true,style=alphabetic,maxbibnames=50]{biblatex}
\addbibresource{/home/nqbh/reference/bib.bib}
\usepackage[utf8]{vietnam}
\usepackage{tocloft}
\renewcommand{\cftsecleader}{\cftdotfill{\cftdotsep}}
\usepackage[colorlinks=true,linkcolor=blue,urlcolor=red,citecolor=magenta]{hyperref}
\usepackage{amsmath,amssymb,amsthm,enumitem,float,graphicx,mathtools,tikz}
\usetikzlibrary{angles,calc,intersections,matrix,patterns,quotes,shadings}
\allowdisplaybreaks
\newtheorem{assumption}{Assumption}
\newtheorem{baitoan}{Bài toán}
\newtheorem{cauhoi}{Câu hỏi}
\newtheorem{conjecture}{Conjecture}
\newtheorem{corollary}{Corollary}
\newtheorem{dangtoan}{Dạng toán}
\newtheorem{definition}{Definition}
\newtheorem{dinhly}{Định lý}
\newtheorem{dinhnghia}{Định nghĩa}
\newtheorem{example}{Example}
\newtheorem{ghichu}{Ghi chú}
\newtheorem{hequa}{Hệ quả}
\newtheorem{hypothesis}{Hypothesis}
\newtheorem{lemma}{Lemma}
\newtheorem{luuy}{Lưu ý}
\newtheorem{nhanxet}{Nhận xét}
\newtheorem{notation}{Notation}
\newtheorem{note}{Note}
\newtheorem{principle}{Principle}
\newtheorem{problem}{Problem}
\newtheorem{proposition}{Proposition}
\newtheorem{question}{Question}
\newtheorem{remark}{Remark}
\newtheorem{theorem}{Theorem}
\newtheorem{vidu}{Ví dụ}
\usepackage[margin=2cm]{geometry}
\def\labelitemii{$\circ$}
\DeclareRobustCommand{\divby}{%
	\mathrel{\vbox{\baselineskip.65ex\lineskiplimit0pt\hbox{.}\hbox{.}\hbox{.}}}%
}
\setlist[itemize]{leftmargin=*}
\setlist[enumerate]{leftmargin=*}

\title{Programming Problems in Mathematical Analysis -- Bài Tập Lập Trình Giải Tích}
\author{Nguyễn Quản Bá Hồng\footnote{A scientist- {\it\&} creative artist wannabe, a mathematics {\it\&} computer science lecturer of Department of Artificial Intelligence {\it\&} Data Science (AIDS), School of Technology (SOT), UMT Trường Đại học Quản lý {\it\&} Công nghệ TP.HCM, Hồ Chí Minh City, Việt Nam.\\E-mail: {\sf nguyenquanbahong@gmail.com} {\it\&} {\sf hong.nguyenquanba@umt.edu.vn}. Website: \url{https://nqbh.github.io/}. GitHub: \url{https://github.com/NQBH}.}}
\date{\today}

\begin{document}
\maketitle
\begin{abstract}
	This text is a part of the series {\it Some Topics in Advanced STEM \& Beyond}:
	
	{\sc url}: \url{https://nqbh.github.io/advanced_STEM/}.
	
	Latest version:
	\begin{itemize}
		\item {\it Programming Problems in Mathematical Analysis -- Bài Tập Lập Trình Giải Tích}.
		
		PDF: {\sc url}: \url{.pdf}.
		
		\TeX: {\sc url}: \url{.tex}.
		\item {\it }.
		
		PDF: {\sc url}: \url{.pdf}.
		
		\TeX: {\sc url}: \url{.tex}.
	\end{itemize}
\end{abstract}
\tableofcontents

%------------------------------------------------------------------------------%

\section{Sequences}

\begin{problem}[\cite{Tao_analysis_1}, Exercise 6.1.1, p. 114]
	Let $\{a_n\}_{n=0}^\infty$ be a sequence of real numbers, such that $a_{n+1} > a_n$, $\forall n\in\mathbb{N}$. Prove that whenever $m,n\in\mathbb{N}$ such that $m > n$, then $a_m > a_n$. We refer to these sequences as \emph{strictly increasing} sequences.
\end{problem}

\begin{proof}[Proof]
	We fix $n\in\mathbb{N}$ and proceed by mathematical induction on $m$, specifically for integers $m > n$. We can formalize this by writing $m = n + k$ for some $k\in\mathbb{Z}_{\ge1}$, and inducting on $k$.
	\begin{itemize}
		\item \emph{Base case $k = 1$}: Here $m = n + 1$. By the hypothesis that the sequence is strictly increasing, we immediately have $a_{n+1} > a_n$.
		\item \emph{Inductive step}: Suppose the claim holds for some $k\in\mathbb{Z}_{\ge1}$, i.e., $a_{n + k} > a_n$. We wish to show it holds for $k + 1$, i.e., $a_{n + k + 1} > a_n$. By the given hypothesis applied to the index $n + k$, we know $a_{n + k + 1} > a_{n + k}$. By our inductive assumption, we have $a_{n + k} > a_n$. By the transitivity of the strict inequality relation on real numbers \cite[Prop. 5.4.7]{Tao_analysis_1}, combining these yields $a_{n + k + 1} > a_n$.
	\end{itemize}	
	This closes the induction, proving that $a_m > a_n$ whenever $m > n$.
\end{proof}

\begin{problem}[\cite{Tao_analysis_1}, Exercise 6.1.2, p. 114]
	Let $\{a_n\}_{n=m}^\infty$ be a sequence of real numbers, and let $L\in\mathbb{R}$. Prove that $\{a_n\}_{n=m}^\infty$ converges to $L$ iff, given any $\varepsilon\in\mathbb{R}_{>0}$, one can find an $N\ge m$ such that $|a_n - L|\le\varepsilon$, $\forall n\ge N$.
\end{problem}

\begin{proof}[Proof]
	This follows directly from unwrapping the definitions provided in the text: By Def. 6.1.5, $\{a_n\}_{n=m}^\infty$ converges to $L$ iff it is \emph{eventually $\varepsilon$-close to $L$}, $\forall\varepsilon\in\mathbb{R}_{>0}$; the sequence is \emph{eventually $\varepsilon$-close to $L$} iff there exists an $N\ge m$ such that the tail sequence $\{a_n\}_{n=N}^\infty$ is \emph{$\varepsilon$-close to $L$}; the sequence $\{a_n\}_{n=N}^\infty$ is \emph{$\varepsilon$-close to $L$} iff $a_n$ is $\varepsilon$-close to $L$, $\forall n\ge N$. By Def. 6.1.2, $a_n$ is $\varepsilon$-close to $L$ iff the distance $d(a_n,L)\le\varepsilon$, i.e., $|a_n - L|\le\varepsilon$.
	
	Stringing these logical equivalences together, we immediately obtain that $\{a_n\}_{n=m}^\infty$ converges to $L$ iff, for every $\varepsilon\in\mathbb{R}_{>0}$, there exists an $N\ge m$ such that $|a_n - L|\le\varepsilon$, $\forall n\ge N$.
\end{proof}

\begin{problem}[\cite{Tao_analysis_1}, Exercise 6.1.3, p. 114]
	Let $\{a_n\}_{n=m}^\infty$ be a sequence of real numbers, let $c\in\mathbb{R}$, and let $m'\ge m$ be an integer. Show that $\{a_n\}_{n=m}^\infty$ converges to $c$ iff $\{a_n\}_{n=m'}^\infty$ converges to $c$.
\end{problem}

\begin{proof}[Proof]
	
\end{proof}

\begin{problem}[\cite{Tao_analysis_1}, Exercise 6.1.4, p. 114]
	
\end{problem}

\begin{proof}[Proof]
	
\end{proof}

\begin{problem}[\cite{Tao_analysis_1}, Exercise 6.1.5, p. 114]
	
\end{problem}

\begin{proof}[Proof]
	
\end{proof}

\begin{problem}[\cite{Tao_analysis_1}, Exercise 6.1.6, p. 114]
	
\end{problem}

\begin{proof}[Proof]
	
\end{proof}

\begin{problem}[\cite{Tao_analysis_1}, Exercise 6.1.7, p. 114]
	
\end{problem}

\begin{proof}[Proof]
	
\end{proof}

\begin{problem}[\cite{Tao_analysis_1}, Exercise 6.1.8, p. 114]
	
\end{problem}

\begin{proof}[Proof]
	
\end{proof}

\begin{problem}[\cite{Tao_analysis_1}, Exercise 6.1.9, p. 114]
	
\end{problem}

\begin{proof}[Proof]
	
\end{proof}

\begin{problem}[\cite{Tao_analysis_1}, Exercise 6.1.10, p. 114]
	
\end{problem}

\begin{proof}[Proof]
	
\end{proof}

\begin{baitoan}
	Viết chương trình {\sf C/C++, Python} để kiểm tra 1 dãy số thực $\{a_n\}_{n=0}^\infty\subset\mathbb{R}$ thỏa \cite[Def. 6.1.3]{Tao_analysis_1}.
\end{baitoan}

\begin{baitoan}
	Cho dãy số thực $\{a_n\}_{n=0}^\infty\subset\mathbb{R}$ được xác định bởi hàm
	\begin{verbatim}
double f(int  n) {
    return <formula of f(n) in {\sf C/C++, Python}>
}
	\end{verbatim}
	Tính $N_\varepsilon$ với $\varepsilon\in\mathbb{R}_{>0}$ bất kỳ bằng cách viết hàm
	\begin{verbatim}
int find_N_epsilon(vector<double> a, double epsilon) {
    // {\sf C/C++, Python} code to find N_\epsilon
}
	\end{verbatim}
\end{baitoan}

\begin{baitoan}
	Viết chương trình {\sf C/C++, Python} để kiểm tra 1 dãy số thực $\{a_n\}_{n=0}^\infty\subset\mathbb{R}$ thỏa \cite[Def. 6.1.5]{Tao_analysis_1}.
\end{baitoan}

\begin{baitoan}
	Viết chương trình {\sf C/C++, Python} mô phỏng lại cách tìm supremum \& infimum của 1 dãy số thực $\{a_n\}_{n=1}^\infty$ được cài đặt bởi hàm sau
	\begin{verbatim}
double a(int m, int n) {
    return a_n := f(n)
}
	\end{verbatim}
	với $f:\mathbb{N}^\star\to\mathbb{R}$ là 1 hàm bất kỳ được viết riêng.
\end{baitoan}

\begin{proof}[Giải]
	
\end{proof}

\begin{baitoan}
	Viết chương trình {\sf C/C++, Python} để xuất ra miền $D_f = {\rm dom}(f)$ \& tập giá trị $R_f = {\rm range}(f) = \{f(x);x\in D_f\}$ của 1 hàm $f:D_f\to\mathbb{R}$ được trỏ vào, sử dụng con trỏ hàm. Lưu trữ \& in ra dãy $\{a_n\}_{n = n_0}^N\subset\mathbb{R}$ cảm sinh bởi hàm $f$ với $n_0,N\in\mathbb{Z}_{\ge1}$, $N\ge n_0$ được nhập vào.
\end{baitoan}

\begin{baitoan}
	Xét các tập $X\subset\mathbb{R}$ có dạng
	\begin{align}
		\label{union of intervals of R}
		&X\left(n_1,n_2,n_3,n_4,\{a_i\}_{i=1}^{n_1},\{b_i\}_{i=1}^{n_1},\{c_i\}_{i=1}^{n_2},\{d_i\}_{i=1}^{n_2},\{e_i\}_{i=1}^{n_3},\{f_i\}_{i=1}^{n_3},\{g_i\}_{i=1}^{n_4},\{h_i\}_{i=1}^{n_4}\right)\nonumber\\
		&\coloneqq\left(\bigcup_{i\in[n_1]} [a_i,b_i]\right)\cup\left(\bigcup_{i\in[n_2]} [c_i,d_i)\right)\cup\left(\bigcup_{i\in[n_3]} (e_i,f_i]\right)\cup\left(\bigcup_{i\in[n_4]} (g_i,h_i)\right).
	\end{align}
	Viết chương trình {\sf C/C++, Python} để tính bao đóng (closure) $\overline{X}$ \& tập biên $\partial X\coloneqq\overline{X}\backslash X$ tất cả các điểm bám dính (adherent point) của $X$ nhưng không thuộc $X$.
\end{baitoan}

\begin{baitoan}
	Viết chương trình {\sf C/C++, Python} để kiểm tra lại \cite[Lem. 9.1.11]{Tao_analysis_1} với $X,Y$ có dạng \eqref{union of intervals of R}.
\end{baitoan}

\begin{baitoan}
	Mở rộng 2 bài toán trước cho phép toán hợp, giao, hiệu, etc.
\end{baitoan}

%------------------------------------------------------------------------------%

\section{Miscellaneous}

%------------------------------------------------------------------------------%

\printbibliography[heading=bibintoc]
	
\end{document}